% GENERATED_BY: translate_oc_core_latex_targets.py
% AI_ASSISTED_TRANSLATION_DRAFT: TRUE
% THEOREM_REVIEW_STATUS: PENDING
% THEOREM_REVIEW_MODE: FORMAL_AI_GATE__LATEX_AWARE_MT
% THEOREM_REVIEW_REVIEWER_ID:
% THEOREM_REVIEW_AUTHORITY_CLASS:
% THEOREM_REVIEWED_AT:
% THEOREM_REVIEW_DOSSIER_REF:
% SOURCE_LANGUAGE: EN
% TARGET_LANGUAGE: RU
% SOURCE_EN_REF: content/cycles/cycles_k0.tex
% ================================================================
% ==== FILE: content/cycles/cycles_k0.tex
% ================================================================

\subsubsection{Циклы на \texorpdfstring{$K_0$}{K_0}}
\label{sec:cycles-k0}

Уровень $K_0$ представляет метаонтологическую основу
Ontology of Continua.
Он не обладает ни геометрией, ни временем, ни допустимым пространством состояний в
смысл $\Omega(K)$, нет потоков $J$, нет порогов $\Theta$ и нет
операторы эволюции.
По этой причине понятие цикла не может быть определено на $K_0$.

\subsubsection{Отсутствие пространства состояний и времени}

Циклы требуют:
\begin{enumerate}
  \item непустое пространство состояний $\Omega(K)$;
  \item понятие времени или параметризация $t\in[0,\tau]$;
  \item поле потока $J$, создающее замкнутые траектории.
\end{enumerate}

Ни одна из этих структур не существует на $K_0$.
Формально:
\[
  \Omega(K_0)=\emptyset, \qquad
  \partial\Omega(K_0) \text{ undefined}, \qquad
  J_{K_0} \text{ undefined},
\]
и не существует параметра $t$, который мог бы индексировать траекторию.

Следовательно, цикл
\[
  C:[0,\tau]\to\Omega(K), \quad C(0)=C(\tau),
\]
невозможно сформулировать.

\subsubsection{Структурная причина}

Уровень $K_0$ — это не континуум, а логическое условие возможности.
для континуумов уровней $K_1$ и выше.
Он не устанавливает никакой динамической или геометрической структуры, на основе которой циклы
могло появиться.
Все понятия, связанные с циклом: длина, стабильность, расстояние до границы,
циклические комплексы, или вклады в непрерывность, строго
неприменимо.

\subsubsection{Значение для более высоких уровней}

Отсутствие циклов на $K_0$ подразумевает:
\begin{itemize}
  \item понятие времени $\tau(K)$ начинается только в $K_1$, где
        сначала определяется непрерывная ось;
  \item первые нетривиальные циклы появляются в $K_1$ как тривиальные или
        вырожденные колебательные структуры одномерного поля;
  \item структурные циклы в полном смысле (стабильность, дистанция до
        $\partial\Omega$, взаимодействие с порогами) возникают только при
        уровни $K_2$ и выше;
  \item рождение циклов само по себе является признаком фазы
        переход $K_0\to K_1$.
\end{itemize}

\subsubsection{Summary}

На $K_0$ циклов не существует.
Структура цикла становится значимой только начиная с $K_1$, где время,
Впервые появляются пространство состояний и потоки.
Соответственно, в файле \texttt{cycles\_k0.tex} записан формальный
\emph{заявление о несуществовании}, требуется для внутренней согласованности
таксономии цикла по всей иерархии континуумов.
